\title{Byte Latent Transformer: Patches Scale Better Than Tokens}

\begin{abstract}
We present the Byte Latent Transformer ({\textsc BLT}), a novel large language model (LLM) architecture operating at the byte level, which is the first to attain parity with traditional tokenization-based LLMs in terms of performance at scale, while offering notable gains in inference speed and robustness. 
{\textsc BLT} processes input as bytes that are grouped into adaptively determined patches, using these patches as the principal computational entities. The patch boundaries are established according to the entropy of subsequent bytes, dedicating increased model resources to regions with higher data complexity. We conduct and report the inaugural scaling analysis of byte-level models under controlled FLOP budgets, scaling to 8B parameters and 4T bytes of training data. Our experiments validate that it is practical to scale models that learn directly from raw bytes and do not rely on predefined vocabularies. Efficiency is enhanced during both training and inference, attributed to the model's ability to dynamically aggregate bytes into larger patches where input is more predictable; this also yields qualitative advances in tasks requiring reasoning and robust generalization to rare cases. On the whole, given a fixed inference cost, {\textsc BLT} outperforms tokenization-based models in scalability by enabling simultaneous increases in both patch granularity and model capacity.
\end{abstract}

\section{Introduction}
\label{section:intro}

We present the Byte Latent Transformer~(\textbf{BLT}), a novel architecture that eliminates the need for tokenization by directly modeling raw byte sequences. To our knowledge, BLT is the first such tokenizer-free system to achieve parity with token-based models at large scales, while also demonstrating notable gains in computational efficiency and resilience (\S\ref{sec:robustness}). Conventional large language models (llms) are predominantly trained in an end-to-end manner, yet still depend on a tokenization stage—a heuristic that segments byte sequences into a fixed vocabulary of tokens. This tokenization step introduces biases into how sequences are compressed, resulting in issues like heightened sensitivity to domains and modalities~\citep{dagangetting}, vulnerability to input perturbations~(\S\ref{sec:robustness}), limited representation of orthographic properties~\citep{edman2024cute}, and persistent multilingual disparities~\citep{liang2023xlm,petrov2024language,limisiewicz2024myte}.

Historically, tokenization has played a crucial role, as training llms directly on raw bytes has been prohibitively expensive at scale owing to the increased sequence lengths~\citep{xue2022byt5}. Some previous studies have attempted to address this issue by adopting more efficient variants of self-attention~\citep{el2020characterbert, clark2022canine} or by introducing architectures that do not rely on attention mechanisms~\citep{wang2024mambabyte}~(\S\ref{section:related_work}). Nevertheless, these solutions are mainly advantageous for training \textit{small models}. When working with large-scale models, the primary computational bottleneck within the Transformer architecture is the substantial feed-forward network layers that must process every individual byte, while the attention mechanism itself contributes far less to the overall computational cost.

For optimized compute allocation, we introduce a dynamic and trainable approach for assembling bytes into \textit{patches}~(\S\ref{section:patching}) along with a novel model architecture that integrates both byte-level and patch-level information. Distinct from traditional tokenization, {\textsc BLT} does not depend on a predetermined patch vocabulary. Instead, any combination of bytes can be transformed into latent patch representations using lightweight, trainable encoder and decoder components. Our findings demonstrate that this approach yields \textit{greater} computational efficiency compared to models based on tokenization.

Tokenization-based language models assign equal computational resources to all tokens, regardless of their individual complexity. This approach prioritizes ease of implementation at the expense of computational efficiency, as token boundaries are dictated by compression algorithms that may not align well with the actual difficulty of generating predictions. A foundational principle of our proposed architecture is to allow models to adjust their computational effort in response to the varying demands of different parts of the input. For instance, predicting the final letters of a word is generally a low-complexity, high-certainty task, in contrast to generating the initial word in a sentence, which often involves more uncertainty and requires greater computational resources—a situation where deploying a large transformer becomes beneficial. In line with this principle, {\textsc BLT}'s structure~(\S\ref{section:architecture}) comprises three transformer blocks: two compact, byte-level \textit{local models} and one expansive, global \textit{latent transformer}~(\autoref{fig:arch_high_level}).

To intelligently partition bytes into patches and allocate computational attention accordingly, {\textsc BLT} leverages next-byte prediction entropy to segment the input, thus assembling contextual byte groupings that maintain consistent information content across patches.

We conduct the first scaling analysis of byte-level models constrained by floating point operation count, extending up to 8B parameters and processing as many as 4T training bytes. Our results demonstrate that it is feasible to train large-scale models directly from bytes, entirely bypassing the need for predefined tokenization schemes. Specifically, {\textsc BLT} achieves flop-controlled training performance on par with Llama 3, yet requires as much as 50\% fewer floating point operations at inference time~(\S\ref{section:scaling}).\footnote{The computational cost is quantified as the total number of Floating Point OPerations (flops) utilized.} 

Moreover, operating directly on raw byte sequences confers notable advantages in modeling data distributions with heavy tails. {\textsc BLT} exhibits increased robustness to input noise relative to conventional tokenizer-driven systems and shows superior capabilities in character-level reasoning, as evidenced in tasks involving orthography, phonological processing, and low-resource machine translation~(\S\ref{sec:robustness}).

Additionally, {\textsc BLT} models allow for concurrent scaling of both the model and the input patch size, all while keeping the inference computational cost constant. Employing longer patches typically leads to computational savings during inference, permitting the allocation of these savings to enlarge the global latent transformer component, since it is invoked less frequently. Through a suite of scaling studies governed by fixed inference flops~(\autoref{fig:fixed_inference_scaling}), we observe substantially more favorable scaling patterns compared to models reliant on tokenization-based encodings.

To summarize, this work offers several key contributions: 1) We present {\textsc BLT}, a novel byte latent llm design that dynamically allocates computational resources to enhance flop efficiency; 2) We demonstrate that our approach matches the flop-controlled training performance of Llama 3 models up to the 8B parameter scale, with the additional ability to exchange small decreases in evaluation metrics for up to a 50\% improvement in flop efficiency; 3) {\textsc BLT} enables a new scaling paradigm for llms, permitting model size increases while keeping inference cost constant; and 4) We provide evidence of {\textsc BLT} models exhibiting superior robustness to input noise, as well as heightened sensitivity to sub-word information that token-based llms typically overlook. All training and inference code for {\textsc BLT} is publicly available at~\url{https://github.com/facebookresearch/blt}.

\section{Patching: From Individual Bytes to Groups of Bytes}
\label{section:patching}

Dividing byte sequences into \textit{patches} enables {\textsc BLT} to flexibly assign computational resources depending on the contextual requirements. As illustrated in Figure~\ref{fig:patching_types}, multiple strategies exist for segmenting bytes into patches. More precisely, a patching function $f_p$ takes a sequence of bytes $\pmb{x}=\{x_i,|i=1,\ldots n\}$ of length $n$ and segments it into a sequence of $m < n$ patches $\pmb{p}=\{p_j|j=1,\ldots,m\}$ by assigning each $x_i$ to either 0 or 1, where a value of 1 denotes the start of a new patch. For models utilizing either token-based or patch-based approaches, the primary determinant of computational effort lies in the number of steps performed by the principal Transformer module. In the context of {\textsc BLT}, this corresponds to the count of patches generated for encoding data according to a given patching function. Therefore, the typical or average patch size, referred to as the \textit{patch size}, serves as the critical metric governing computational requirements for both training and inference under a selected patching scheme~(\S\ref{section:flops}).  
We proceed by outlining three specific patching methods: fixed-length patching~(\S\ref{section:static-patch}), whitespace-driven patching~(\S\ref{section:white-patch}), and dynamic patching guided by byte-level language model entropies~(\S\ref{section:dyn-patch}). We conclude by considering incremental patching techniques and clarifying the distinctions between patching and traditional tokenization~(\S\ref{section:bpe-patch}).

\subsection{Strided Patching Every K Bytes}
\label{section:static-patch}

A common and simple method for partitioning bytes involves creating fixed-size patches of $k$, as implemented in MegaByte~\citep{yu2023megabyte}. This approach, utilizing a constant stride, is straightforward to use during both training and inference. It also allows for an uncomplicated adjustment of the mean patch size, granting easy manipulation of the computational cost in terms of floating point operations. Despite these advantages, the method has notable limitations. Most importantly, computational resources are statically assigned, regardless of the demands of specific regions: computational steps may be wasted when handling uninformative tokens such as whitespace in code, or, conversely, may be insufficient for byte segments rich in content like mathematical notation. Furthermore, this technique can result in the uneven and context-insensitive splitting of byte sequences, leading to inconsistencies where identical words are divided in varying ways across occurrences.

\subsection{Space Patching}
\label{section:white-patch}

\citet{slagle2024spacebyte} introduces a straightforward but impactful refinement to strided patching by forming new patches immediately following any space-like bytes\footnote{
    Space-like bytes refer to any byte that is neither a Latin letter, digit, nor utf-8 continuation byte. Additionally, each patch is required to include at least one byte that is not space-like. }. These bytes frequently correspond to natural linguistic boundaries across many languages. Under the Space patching approach, each word is modeled with an additional latent transformer computation step (i.e., increased flops), guaranteeing uniform patching for words within different sequences. This allocation of computational resources targets predictions that are typically more challenging and tend to occur after spaces. To illustrate, the prediction of the initial byte in the response to the prompt ``Who composed the Magic Flute?\uline{\hspace{.6em}}'' presents greater difficulty than predicting the subsequent bytes after ``M'', as that initial byte substantially narrows down the set of plausible completions, thereby making the full answer ``Mozart'' much easier to model thereafter. Nonetheless, the Space patching method struggles with coverage across all languages and domains, and crucially, does not support flexible patch sizing. In the following section, we propose a novel patching technique that leverages the observation that word-initial bytes pose the highest predictive difficulty, while simultaneously enabling straightforward control over patch sizes.

\subsection{Entropy Patching: Using Next-Byte Entropies from a Small Byte LM}
\label{section:dyn-patch}

Instead of employing traditional rule-based heuristics like whitespace, our methodology leverages a data-centric framework to pinpoint next-byte predictions that exhibit high uncertainty. We present \textit{entropy patching}, a technique that determines patch boundaries through entropy estimations.

To accomplish this, we fit a compact auto-regressive language model operating at the byte level on the {\textsc BLT} training dataset, and then calculate the entropies of the next-byte distributions $p_{e}$ over the vocabulary $\mathcal{V}$ as determined by the language model:

\begin{equation}
H(x_i) = \sum_{v \in \mathcal{V}} p_{e}(x_i=v|\pmb{x}_{<i}) \log p_{e}(x_i=v|\pmb{x}_{<i})
\end{equation}

We investigate two approaches for detecting patch boundaries using entropies $H(x_i)$. The first method selects points that exceed a predefined global entropy threshold, as shown in~\autoref{fig:patching}. The second strategy focuses on locating points with entropy values substantially higher than the preceding ones. This latter technique may also be viewed as pinpointing instances where the entropy, which tends to decrease monotonically within a patch, experiences a local increase.

\begin{align*}
    \text{Global Constraint}\hspace{1cm}H(x_t)& > \theta_g\\
    \text{Approx. Monotonic Constraint}\hspace{1cm}H(x_t)& - H(x_{t-1}) > \theta_r
\end{align*}

Patch boundaries are determined through a minimal preprocessing procedure that takes place as part of the dataloading phase. This approach contrasts with the method used by~\citet{nawrot-etal-2023-efficient}, who train a classifier specifically to identify entropy-based patch boundaries. Within our experiments~(\S\ref{section:experiments}), we perform a comparative analysis of these two techniques for classifying bytes as either low or high entropy.

\subsection{The Byte-Pair Encoding (BPE) Tokenizer and Incremental Patching}
\label{section:incremental-patching}
\label{section:bpe-patch}

A significant number of contemporary llms, such as our reference model Llama 3, rely on subword tokenizers, notably bpe~\citep{gage1994new, sennrich-etal-2016-neural}. In this context, we define ``tokens'' as groups of bytes selected from a \textit{finite} vocabulary that is established before training begins. In contrast, ``patches'' denote adaptively created sequences that are not restricted by a predefined vocabulary. One important distinction between tokens and patches lies in the model’s accessibility: when using tokens, the underlying byte-level information remains inaccessible to the model.

An important advancement introduced by {\textsc BLT} compared to tokenization-based models lies in its redefinition of the relationship between vocabulary size and computational demands. In conventional large language models, enlarging the vocabulary generally leads to larger average token sizes, which reduces the number of inference steps but simultaneously results in a greater output dimensionality in the model's final projection layer. This dynamic constrains tokenization-based methods, limiting their capacity to significantly vary token sizes and inference efficiency. As an illustration, Llama 3 raises the average token size from 3.7 to 4.4 bytes, but this comes at the expense of expanding its embedding table by a factor of 4 relative to Llama 2.

During the generation process, {\textsc BLT} must determine whether the current position within the byte sequence aligns with a patch boundary, as this influences whether additional computation through the Latent Transformer is required. Importantly, this decision must be made based solely on the already-generated portion of the sequence, without reference to bytes that have not yet been produced. As a result, patch segmentation cannot rely on future bytes to inform its decisions. Mathematically, a patching function $f_p$ is considered incrementally patching if the following holds:
$$f_p(\pmb{x}_{<i}) = f_p(\pmb{x})_{<i}$$
The bpe algorithm does not qualify as an incremental patching method, as it can segment the same initial subsequence in various ways depending on the continuation that follows, and therefore does not adhere to the stated property.\footnote{Introducing a dedicated delimiter token to mark patch boundaries allows bpe to behave as an incremental patching method, but at the cost of lengthening the overall byte sequence.} 

\section{{\textsc BLT} Architecture}
\label{section:architecture}
The {\textsc BLT} framework consists of a sizeable, global autoregressive language model that processes patch-level representations. Complementing this core model are two smaller, localized modules tasked with mapping byte sequences into patch representations and reconstructing the original byte sequences from these patches~(\autoref{fig:arch_high_level}).

\subsection{Latent Global Transformer Model}

\textit{The Latent Global Transformer} refers to an autoregressive transformer model $\mathcal{G}$ comprised of $l_{\mathcal{G}}$ layers. This model transforms a sequence of latent input patch embeddings, $p_j$, into corresponding output patch embeddings, $o_j$. In our notation, subscripts $j$ and $i$ are reserved for identifying patches and bytes, respectively. The architecture employs a block-causal attention mask as introduced by~\citep{dubey2024llama}, which limits attention to the current and preceding patches within the same document. The global model is responsible for the majority of computational cost, both in pre-training and inference. Consequently, selectively invoking this component enables flexible allocation of computational resources depending on the complexity associated with different sections of the input and output sequences.

\subsection{Local Encoder}
\label{section:local}

\textit{The Local Encoder Model}, represented as $\mathcal{E}$, is designed as a compact transformer architecture comprising $l_{\mathcal{E}}$ layers, where $l_{\mathcal{E}}$ is much smaller than $l_{\mathcal{G}}$. Its principal function is to rapidly transform a series of input bytes $b_i$ into rich patch-level representations $p_j$. Unlike standard transformers, this model incorporates a cross-attention mechanism after each transformer block, allowing it to aggregate the byte-level embeddings into patch representations~(\autoref{fig:crossattn}). 

Initially, the sequence of input bytes $b_i$ is projected via a learned embedding matrix in $\mathbb{R}^{256 \times h_{\mathcal{E}}}$, resulting in token embeddings $x_i$. These vectors can be further enriched using hash-based embeddings when appropriate~(\S\ref{sec:hash-emb}). 

Subsequently, the representation passes through stacked layers of transformers and cross-attention modules~(\S\ref{sec:enc-cross-attn}), producing patch-level embeddings $p_i$ that serve as input to the downstream global transformer, $\mathcal{G}$. The transformer modules employ a \textit{local block causal} attention masking scheme, where each byte token is allowed to focus on a sliding window of $w_{\mathcal{E}}$ prior bytes. This attention mechanism may span across patch demarcations but never extends beyond document boundaries. The following subsections provide further elaboration on the embedding strategy and the structure of the cross-attention module.

\subsubsection{Encoder Hash n-gram Embeddings}
\label{sec:ngram-emb}
\label{sec:hash-emb}
An essential aspect of building reliable and rich representations at every position $i$ involves integrating knowledge of the prior bytes. In {\textsc BLT}, this is accomplished by capturing information about each byte $b_i$ on its own as well as within the context of a byte n-gram. Specifically, at every index $i$, we begin by generating byte-grams

\begin{align}
    g_{i,n}=\{b_{i-n + 1},\ldots, b_i\}
\end{align}

for each byte position $i$ and for values of $n$ ranging from three to eight.\footnote{
Byte-grams with length $n$ or greater are excluded if $i < n$.}

Next, we propose hash $n$-gram embeddings, in which all byte $n$-grams are projected through a hash function to a position in a dedicated embedding matrix $E_{n}^{hash}$, each matrix having a predetermined dimensionality for every $n \in\{3,4,5,6,7,8\}$~\citep{bai2010learning}.
The embedding produced in this way is combined with the embedding of the byte itself, after which the sum is normalized and used as the input to the local encoder model. The enriched embedding is computed as follows:

\begin{align}
e_i &= x_i + \sum_{n = 3,...,8}E_{n}^{hash}(\text{Hash}(g_{i,n})) \\
\text{where, Hash}(g_{i,n}) &= \text{RollPolyHash}(g_{i,n}) \% |E_{n}^{hash}|
\end{align}

We adjust $e_i$ by dividing by the total count of $n$-gram sizes incremented by one, and employ the RollPolyHash method as outlined in Appendix~\ref{appendix:rollpolyhash}. Section~\ref{section:ablations} presents an ablation study analyzing how varying $n$-gram hash embedding parameters, specifically the value of $n$ and the size of the embedding table, influence the scaling law behavior under controlled FLOP budgets. Beyond utilizing hashed $n$-gram embeddings, we conducted experiments with $n$-gram embeddings based on frequency, for which we detail our methodology and findings in Appendix~\ref{appendix:ngrams}.

\subsubsection{Encoder Multi-Headed Cross-Attention}
\label{sec:enc-cross-attn}
Our approach largely mirrors the cross-attention input layer utilized in the Perceiver architecture~\citep{jaegle2021perceiver}, with the key distinction that here, each latent representation is associated with an individual, variable-length patch representation rather than a predetermined set of latent variables~(\autoref{fig:crossattn}), and focuses attention exclusively on the bytes forming the given patch. This component uses a query vector for every patch $p_j$, initialized by aggregating (pooling) the byte-level embeddings belonging to that patch, after which a linear projection, $\mathcal{E}_{C} \in \mathbb{R}^{h_{\mathcal{E}} \times (h_{\mathcal{E}} \times U_{\mathcal{E}})}$, is applied. Here, $U_{\mathcal{E}}$ represents the count of cross-attention heads within the encoder. More precisely, denoting the sequence of bytes in patch $p_j$ as $f_{\text{bytes}}(p_j)$, we compute

\begin{align}
P_{0,j} &= \mathcal{E}_{C}(f_\text{bytes}((p_j)), f~ \text{is a pooling function} \\
P_l &= P_{l-1} + W_o\left(\text{softmax}\left(\frac{QK^T}{\sqrt{d_k}}\right)V\right) \\
\text{where } Q_j &= W_q(P_{l-1,j}), K_i = W_k(h_{l-1,i}), V_i = W_v(h_{l-1,i}) \\
h_l &= \text{Encoder-Transformer-Layer}_l(h_{l-1})
\end{align}

Here, $P \in \mathbb{R}^{n_p \times h_{\mathcal{G}}}$ denotes the set of $n_p$ patch representations to be fed into the global model, each of which is initially formed by aggregating the byte embeddings $e_i$ that are associated with the bytes of patch $p_j$. The matrices $W_q$, $W_k$, $W_v$, and $W_o$ serve as the parameterizations for query, key, value, and output projections, respectively, where both keys and values are derived from projecting the byte representations $h_i$ of the previous layer (or $e_i$ for the initial layer). A tailored masking method is implemented so that, for each patch, the corresponding query $Q_j$ is restricted to attend only to those keys and values pertaining to the bytes of patch $j$. Given that cross-attention utilizes a multi-head attention mechanism over $Q$, $K$, and $V$, and since patch representations generally have a higher dimensionality ($h_{\mathcal{G}}$) relative to $h_{\mathcal{E}}$, we organize $P_l$ as multiple attention heads of dimension $h_{\mathcal{E}}$ during cross-attention and subsequently concatenate these heads to recover the $h_{\mathcal{G}}$-dimensional space. Furthermore, pre-LayerNorm is applied to queries, keys, and values within this module, and positional encodings are omitted. The architecture incorporates a residual pathway encircling the cross-attention block.

\subsection{Local Decoder} 

The local decoder $\mathcal{D}$, analogous to the local encoder, is designed as a compact transformer-based architecture comprising $l_{\mathcal{D}} << l_{\mathcal{G}}$ layers. Its primary function is to translate a sequence of global patch embeddings $o_j$ into a sequence of raw bytes, $y_i$. At each step, the local decoder generates the next byte by conditioning on all previously generated bytes, utilizing the hidden states output by the local encoder for the byte sequence as input. Within the local decoder, there is a repeated stack of $l_{\mathcal{D}}$ layers, each consisting of a cross-attention module followed by a transformer block. Notably, the cross-attention mechanism precedes the transformer block in each layer, allowing the decoder to derive byte-level representations from the provided patch representations. Subsequent processing of this byte sequence is then carried out by the transformer's local decoder layers.

\subsubsection{Decoder Multi-headed Cross-Attention} 

Within the decoder's cross-attention mechanism, the queries consist of the byte-representations, while the key and value vectors correspond to the patch representations, effectively reversing their typical assignments. For the cross-attention, the starting byte-representations are derived from the byte embeddings output by the final encoder layer, specifically $h_{l_{\mathcal{E}}}$. For each layer $l$, the updated byte-representations $d_{l,i}$ are then calculated as follows:

\begin{align}
D_0 &= h_{l_{\mathcal{E}}} \\
B_l &= D_{l-1} + W_o\left(\text{softmax}\left(\frac{QK^T}{\sqrt{d_k}}\right)V\right), \\
  &\text{where } Q_i = W_q(d_{l-1,i}), K_i = W_k(\mathcal{D}_C(o_j)), V_i = W_v(\mathcal{D}_C(o_j)) \\
  D_l &= \text{Decoder-Transformer-layer}_l(B_l)
\end{align}

In this formulation, $W_k$ and $W_v$ denote the key and value projection matrices, respectively, each carrying out a linear transformation followed by a split operation $\mathcal{D}_C$ on the ultimate patch representations $o_j$ generated by the global model. $W_q$ is designated as the query projection matrix and acts upon byte representations $d_{l-1}$ from the preceding decoder transformer layer (or $h_{l_{\mathcal{E}}}$ in the case of the initial layer). The matrix $W_o$ serves as the output projection, such that $B \in \mathbb{R}^{h_{\mathcal{D}} \times n_b}$, where $n_b$ represents the count of output bytes. The subsequent decoder representations, $D_l$, are derived by applying a decoder transformer layer to the result of the cross-attention block, $B$. Mirroring the methodology used in the local encoder cross-attention, our approach involves the deployment of multiple attention heads, employs pre LayerNorm, omits positional embeddings, and incorporates a residual connection enveloping the cross-attention module.

\section{Experimental Setup}
\label{section:experiments}
We meticulously craft controlled experiments to evaluate {\textsc BLT} alongside tokenization-based models, ensuring that {\textsc BLT} is not afforded any undue benefit from potentially accessing longer sequence contexts.

\subsection{Pre-training Datasets} The models discussed in this work are pre-trained on two main datasets: 1) The Llama 2 dataset~\citep{touvron2023llama}, containing 2 trillion tokens sourced from diverse public data, which undergo extensive cleaning and filtering processes to ensure high data quality; and 2) {\textsc BLT}-1T: a newly assembled corpus comprising 1 trillion tokens, also drawn from a range of publicly available sources and incorporating a selection of the pre-training data made available by Datacomp-LM~\citep{li2024datacomp}. The Llama 2 dataset serves as the basis for scaling law analyses focused on identifying the optimal token count, in accordance with findings by~\cite{dubey2024llama}, thereby guiding architectural decisions for {\textsc BLT}. In contrast, {\textsc BLT}-1T is employed for full-scale pre-training to facilitate direct performance comparisons with Llama 3 on various downstream benchmarks. Importantly, neither dataset includes any information from Meta platforms or services. Additionally, for models utilizing a tokenizer, we employ the Llama 3 tokenizer, which has a 128K vocabulary size, as it delivered superior baseline results compared to the Llama 2 tokenizer in our experimental evaluations.

\subsection{Entropy Model} For our experiments, the entropy model is implemented as a byte-level language model, trained on the identical dataset as the main {\textsc BLT} model. Except where stated otherwise, the architecture consists of a transformer comprising 100M parameters distributed over 14 layers, with each layer having a hidden size of 512. The attention mechanism operates over a sliding window of 512 bytes. All other hyperparameters mirror those utilized in the local and global transformer configurations. We also explored variations in model scale, the size of the receptive field, and alternative architectures, as detailed in \autoref{section:ablations}. Notably, when the receptive field is sufficiently constrained, it becomes feasible to represent the trained entropy model as a compact and efficient lookup table.

\subsection{Entropy Threshold and Equalizing Context Length} For models employing entropy-driven patching, we determine a patching threshold that yields a specified average \textit{patch size} across the pretraining data mixture. In the case of {\textsc BLT}, as opposed to traditional tokenization, the \textit{patch size} can be selected freely, which can greatly affect the effective context length available to the model. To ensure comparability and prevent models with larger patch sizes from benefiting unduly from longer context, we fix the expected total number of bytes per batch. Consequently, for configurations utilizing larger patch sizes, we proportionally decrease the sequence length. Specifically, on Llama 2 data, the context is set to 8k bytes, whereas for the {\textsc BLT}-1T dataset, the average context window is expanded to 16k bytes—always holding the batch size at an average of 16M bytes.

Although the mean batch size remains fixed, employing dynamic patching techniques results in varying proportions between the number of bytes and patches within each batch. To enhance efficiency, our {\textsc BLT} training pipeline organizes patches into batches in a manner that eliminates the need for padding operations within the computationally intensive latent transformer component. This strategy guarantees a uniform patch count per batch. To further control memory usage and prevent surges caused by exceptionally large patches, we apply padding—and, when necessary, truncation—to the byte sequences, capping them at 12k bytes for Llama 2 and 24k bytes for the {\textsc BLT}-1T datasets during training.

\subsection{Entropy Model Context}
\label{section:entropy-context}
Our empirical observations reveal that applying entropy patching results in increasingly large patches, particularly in structured formats such as multiple choice tasks, which characteristically exhibit a high degree of repetition (see an example involving MMLU in \autoref{fig:mmlu_patching}). These differences arise from the reduced entropy associated with repeated content within the entropy model context. Accordingly, for the primary large-scale experiment utilizing {\textsc BLT}-Entropy with a patch size of 4.5, we refresh the entropy context by introducing new lines and employ an approximate monotonicity constraint, as this approach demonstrates greater resilience against “entropy drift” due to varying context length. While this adjustment modifies our entropy computation method, the overall approach for determining the entropy threshold remains unchanged.

\subsection{FLOPs Estimation} 
\label{section:flops}

Our approach closely mirrors the calculation of transformer FLOPs outlined in Chinchilla~\citep{hoffmann2022training}, accounting for the FLOPs associated with the feed-forward modules, the qkvo projections within the self-attention block, and the computations for attention as well as the output projection. However, we diverge by modeling the input embedding layer as an efficient lookup rather than performing a dense matrix multiplication, effectively rendering this operation flop-free. Consistent with prior research, we approximate that the backward pass entails double the number of FLOPs compared to the forward pass.

In order to determine the number of flops \textit{per byte} for {\textsc BLT} models, we aggregate the flops contributed by the local encoder transformer, the global latent transformer, and the local decoder transformer, as well as those from the cross-attention components within both the encoder and decoder:

\begin{align}
    \text{FL}_{\text{{\textsc BLT}}} &= \text{Transf. FL}(h_{\mathcal{G}}, l_{\mathcal{G}}, m=n_{ctx}/n_p,V=0) / n_p\\
   &+\text{Transf. FL}(h_{\mathcal{E}}, l_{\mathcal{E}}, m=w_{\mathcal{E}}, V=0) \\
   &+ \text{Transf. FL}(h_{\mathcal{D}}, l_{\mathcal{D}}, m=w_{\mathcal{D}}, V=256) \\ 
    &+ \text{Cross Attn. FL}(h_{\mathcal{E}}, l_{\mathcal{E}}, m=n_p, r=n_p/k) \times k/n_p \\ 
   &+ \text{Cross Attn. FL}(h_{\mathcal{D}}, l_{\mathcal{D}}, m=k, r=k/n_p) 
\end{align}

In this context, $n_{ctx}$ denotes the length of the sequence measured in bytes, while $n_p$ indicates the patch size. The parameter $r$ specifies the proportion of queries relative to keys and values, and $k$ represents the factor by which the patch dimension exceeds the byte dimension; in other words, $k$ indicates how many local model segments are concatenated to produce the full model representation ($k=2$ as illustrated in \autoref{fig:crossattn}). $V$ stands for the vocabulary size associated with the output projection layer, which is utilized exclusively within the local decoder. The length of the attention context, $m$, varies according to whether the operation is performed on the byte sequence or the patch sequence, prompting an adjustment in the computation of attention flops for each module. The precise formulations for the Transformer-FLOPs and Cross-Attention FLOPs can be found in Appendix~\ref{section:flops-appendix}.

\subsection{Bits-Per-Byte Estimation} Perplexity is meaningful only when a tokenizer is fixed, since it quantifies token-level prediction uncertainty. However, for fair comparisons between models operating at the byte and token levels, as established in prior studies~\citep{xue2022byt5, yu2023megabyte, wang2024mambabyte}, we utilize the Bits-Per-Byte (BPB) metric. BPB serves as a perplexity-like measure that does not depend on any particular tokenization scheme. Formally:

\begin{align}
    \text{BPB}(x)&= \frac{\mathcal{L}_{CE}(\pmb{x})}{\ln(2)\cdot n_{\text{bytes}}}
\end{align}

Here, the total cross-entropy loss for the dataset $\pmb{x}$ is divided by the product of the total number of bytes in $\pmb{x}$ and the conversion factor $\ln(2)$, thus standardizing model uncertainty across bytes irrespective of the tokenizer.

\subsection{Transformer Architecture Hyperparameters} 
For every transformer block in {\textsc BLT}—encompassing both the local and global model variants—we closely adhere to the Llama~3~\citep{dubey2024llama} design choices. Specifically, we employ the SwiGLU activation function~\citep{swiglu} within the feed-forward networks, integrate rotary positional embeddings (RoPE)~\citep{RoPe} with a parameter setting of $\theta=500000$~\citep{xiong2024effective} exclusively in self-attention modules, and apply RMSNorm \citep{RMSNorm} for the normalization of layers. For self-attention operations utilizing fixed-standard attention masks—such as \textit{block causal} or \textit{fixed-window block causal}—Flash Attention~\citep{dao2022flashattention} is utilized, adopting a window size of 512 for masks with fixed width. Due to the cross-attention layers requiring dynamically generated, patch-dependent attention masks, we leverage Flex Attention\footnote{\url{https://pytorch.org/blog/flexattention}} for producing fused kernel implementations, leading to substantial acceleration in the training process.

\subsection{{\textsc BLT}-Specific Hyperparameters} 

To evaluate the performance of {\textsc BLT} models, our experiments focus on two main areas: examining scaling behavior and assessing downstream task performance. For these analyses, we experiment with model sizes of 400M, 1B, 2B, 4B, and 8B parameters. Detailed architectural hyperparameters for these configurations are included in Appendix~Table~\ref{tab:arch}.

Query vectors for the initial cross-attention layer within the local encoder are initialized using max-pooling. Each {\textsc BLT} model employs $500,000$ hashes and a single hash function, utilizing n-grams of length 3 to 8. Across all model sizes, a consistent learning rate of $4e-4$ is applied. This learning rate was determined as optimal based on a hyperparameter search between $1e-3$ and $1e-4$ conducted at the 400M and 1B model scales, revealing no difference in optimal learning rates between token and {\textsc BLT} variants.

For scaling analyses with Llama-2 data, training batch sizes follow the guidelines specified in~\cite{dubey2024llama} or are adjusted to their byte-equivalent. Model optimization is performed using the AdamW optimizer~\citep{adamw}, configured with $\beta_1 = 0.9$, $\beta_2 = 0.95$, and $\epsilon = 10^{-8}$. A learning rate schedule comprising a linear warm-up over 2000 steps, followed by cosine decay to zero, is employed. Additionally, we use a weight decay coefficient of 0.1 and enforce global gradient clipping with a threshold set at 1.0.

\section{Scaling Trends}
\label{section:scaling}

We offer a comprehensive overview of the scaling characteristics associated with byte-level models, which provides valuable guidance for the ongoing expansion of {\textsc BLT} models. Our investigation into scaling is designed to overcome the shortcomings identified in earlier studies of byte-level models through the following approaches: (a) We analyze scaling behavior within compute-optimal training paradigms, (b) We train corresponding 8B parameter models on substantial datasets (reaching up to 1T tokens or 4T bytes) and conduct evaluations using downstream benchmarks, and (c) We assess scaling trends while keeping inference cost fixed. A subsequent section will delve into the particular benefits of modeling data as byte sequences.

\subsection{Parameter Matched Compute Optimal Scaling Trends}
\label{section:parameter-matched-scaling}
Leveraging the Llama 2 dataset, we conduct training on both \textit{compute-optimal} bpe and {\textsc BLT} model architectures, spanning four distinct scales from 1B up to 8B parameters. Subsequently, we visualize the relationship between training flops and language modeling efficacy on a representative sample from the data mixture used for training. The bpe models utilize the theoretically optimal allocation between model size and training data, as specified by the guidelines from Llama~3~\citep{dubey2024llama}. This \textit{compute-optimal} configuration is postulated to maximize model performance under a constrained training compute budget~\citep{hoffmann2022training}, establishing a rigorous benchmark for evaluation. For each bpe model in our experiments, we train an equivalent {\textsc BLT} counterpart with an identical Latent Transformer size and architecture, utilizing the same dataset to ensure comparability.

As shown in~\autoref{fig:scaling-compute-opt} (right), {\textsc BLT} models consistently equal or surpass the performance of their bpe counterparts, with this pattern persisting across increases in both model size and computational resources. To our knowledge, {\textsc BLT} represents the first byte-level Transformer architecture to exhibit scaling behaviors on par with BPE-based models within compute-optimal settings. This observation supports our hypothesis that the optimal parameter-to-compute ratio established for bpe models is also applicable to {\textsc BLT}, or at least is closely aligned.

Achieving competitive bpe scaling trends necessitates both advances in architecture and the use of dynamic patching. As shown in ~\autoref{fig:scaling-compute-opt} (left), we evaluate models utilizing space-patching approaches relative to Llama 3. For our assessment of SpaceByte \citep{slagle2024spacebyte}, we replicate it by configuring {\textsc BLT} with space-patching, omitting n-gram embeddings and cross-attention. Despite SpaceByte surpassing Megabyte in performance, there remains a considerable gap between its results and those of Llama 3. In \autoref{fig:scaling-compute-opt} (right), we depict the benefits that result from integrating both enhanced architectural design and dynamic patching. When scaled, {\textsc BLT} models achieve results on par with leading tokenizer-based systems, including Llama 3.

Additionally, our analysis reveals that the particular tokenizer chosen impacts the performance of tokenizer-based models: those utilizing the Llama-3 tokenizer achieve superior results compared to models trained with the Llama-2 tokenizer under otherwise identical data conditions.

Lastly, our {\textsc BLT} architecture exhibits a performance trajectory situated between that of Llama 2 and Llama 3 when utilizing substantially increased patch sizes. While the bpe tokenizers used in Llama 2 and 3 produce tokens averaging 3.7 and 4.4 bytes respectively, {\textsc BLT} demonstrates comparable scaling dynamics with average patch sizes reaching 6 and even 8 bytes. Since inference flop count is inversely related to the average patch size, employing a patch size of 8 bytes can result in inference flop reductions of nearly 50\%. Moreover, models employing larger patch sizes tend to yield enhanced performance as both model and data scales are increased. For example, although {\textsc BLT} with an 8-byte patch size performs markedly worse than bpe Llama 2 at the 1B parameter scale, it surpasses the bpe baseline by the 7B scale. This observation points to the possibility that such patch sizes may yield superior results as models and training datasets continue to scale, and it raises the prospect that even larger patch sizes could become practical as model capacity and compute resources increase further.

\subsection{Beyond Compute Optimal Task Evaluations}
\label{section:evals}
To more thoroughly investigate scaling characteristics, we train an 8B {\textsc BLT} model at a parameter-to-data ratio exceeding the compute optimal on the {\textsc BLT}-1T dataset—a larger and higher-quality collection—and assess its capabilities across a range of widely used classification and generation benchmarks. The evaluation spans a set of tasks that probe common sense reasoning, general knowledge, and code synthesis:
\paragraph{Classification tasks} consist of ARC-Easy (0-shot)~\citep{Eval_ARC}, Arc-Challenge (0-shot)~\citep{Eval_ARC}, HellaSwag (0-shot)~\citep{Eval_hellaswag}, PIQA (0-shot)~\citep{Eval_piqa}, and MMLU (5-shot)~\citep{hendrycks2020measuring}. We utilize a prompt-based scoring procedure by computing the model-assigned likelihoods to the choice options, and summarize performance using mean accuracy.
\paragraph{Coding related generation tasks:} We evaluate code generation capabilities of the model via pass@1 on MBPP (3-shot)~\citep{Eval_MBPP} and HumanEval (0-shot)~\citep{Eval_HumanEval}, focusing on Python code synthesis proficiency of large language models.

In \autoref{tab:evals}, we present a comparison of three models, each trained on the {\textsc BLT}-1T dataset: a model using a bpe Llama 3 tokenizer,\footnote{We select the Llama 3 tokenizer, which has a 128k vocabulary, as it yields superior results compared to the 32k vocabulary of Llama 2.} and two distinct {\textsc BLT} variants. The first variant applies a space-patching strategy ({\textsc BLT}-Space), while the second utilizes a patching strategy based on entropy ({\textsc BLT}-Entropy). The latter incorporates an approximate monotonicity constraint and resets the entropy model's context with new lines, following the method described in~\autoref{section:entropy-context}. All three models are trained under the same computational (flop) budget. For {\textsc BLT}-Entropy, we also experiment with lowering the inference entropy threshold from 0.6 to 0.1, which results in improved task performance at the expense of requiring more inference steps.

The {\textsc BLT}-Entropy model achieves superior results compared to the Llama 3 model on 4 of the 7 evaluated tasks, despite both models being exposed to an equivalent amount of training data. This enhanced performance likely arises from two main factors: (1) more efficient utilization of training computation enabled by dynamic patching, and (2) the explicit modeling of data at the byte level rather than relying on token-based representations.

In contrast, {\textsc BLT}-Space lags behind the Llama 3 tokenizer across all tasks except one; however, it offers a notable decrease in inference flops owing to its greater average patch size of 6 bytes. By way of comparison, both the bpe and entropy-patching approaches yield similar average patch sizes, close to 4.5 bytes, when evaluated on the combined training dataset. Given an equivalent training budget, the model utilizing the larger patch size is able to process 30\% more data than the bpe and entropy-patching models, which may in turn cause {\textsc BLT} to deviate further from the point of compute optimality.

\subsection{Patches Scale Better Than Tokens} {\textsc BLT} models allow for concurrent scaling of both the model and patch dimensions, all while holding the training and inference flop budget, as well as the training dataset size, constant. Unlike token-based models with a fixed vocabulary, patch-based approaches offer the distinctive advantage of enabling unrestricted increases in patch size, circumventing the efficiency limitations described in Section~\ref{section:bpe-patch}. Utilizing larger patches leads to reduced computational requirements, thereby freeing resources that can be redirected to increase the capacity of the global latent transformer, since its application frequency is diminished.

We perform a fixed inference scaling analysis to evaluate whether larger models executing fewer iterations on larger patches can outperform smaller models that require more steps. Utilizing models with 400m and 3.6B parameters initialized with the Llama 2 tokenizer, we identify flop-equivalent models configured with the Llama 3 tokenizer as well as {\textsc BLT}-Entropy models that employ average patch sizes of 6 and 8 bytes over the training datamix (refer to~\autoref{tab:fixed-inf-params} for model specifications). For those models employing a patch size of 8, three encoder layers are utilized in place of one. Each model is trained across different training flop budget allocations.

\autoref{fig:fixed_inference_scaling} illustrates that {\textsc BLT} models display superior scaling behavior compared to tokenization-based models across both categories of inference flops. Initially, BPE models demonstrate stronger performance when the training budget is limited, yet they are soon outperformed by {\textsc BLT} models once the budget exceeds the compute-optimal threshold. From a practical standpoint, allocating additional resources for the one-time pretraining phase may be advantageous, as it can yield a model with improved performance under a constant inference cost. A notable case in point is the family of 8B models, such as Llama 3.1, which has been trained on data volumes two orders of magnitude larger than what would be considered optimal for its parameter scale.

The point at which {\textsc BLT} surpasses token-based architectures has shifted marginally closer to the compute-optimal regime as model sizes increase (from being 3 times to now approximately 2.5 times the compute-optimal budget for larger flop classes). In the same vein, models utilizing a larger patch size of 8 exhibit more rapid scaling performance within higher flop classes, surpassing smaller patch counterparts earlier in the scaling process. As highlighted in Section~\ref{section:parameter-matched-scaling}, increasing patch sizes tend to align {\textsc BLT}'s performance more closely with that of BPE models when operating at larger model scales. We partially ascribe this improvement to the relatively smaller fraction of total computational cost attributed to the byte-level Encoder and Decoder modules—modules that scale at a slower rate than the Latent Transformer. When scaling up total parameters twentyfold (from 400M to 8B), the parameter count specific to the local modules within {\textsc BLT} increases only by about a factor of two. This detail is significant because larger patch sizes influence only the computational demand of the patch Latent Transformer and leave the byte-level modules unaffected. For this reason, the model size ratio for {\textsc BLT}-Entropy with ps=8 increased only slightly, from 1.6x to 1.7x relative to the Llama 2 model size when transitioning to a larger model scale.

In conclusion, our investigation into patch-length scaling reveals that the {\textsc BLT} patch-based model architecture attains more favorable scaling behaviors when both the patch size and the model capacity are increased concurrently. These advantageous scaling characteristics not only remain stable but also appear to strengthen as model size grows further.

\section{Byte Modeling Improves Robustness} We additionally assess the robustness of {\textsc BLT} in contrast to token-based models without inherent byte-level representation, and introduce a technique for augmenting pretrained token-based models with byte-level information.  
\label{sec:robustness}

\subsection{Character-Level Tasks}

The initial impetus for developing byte-level models stemmed from their inherent resilience to noise at the byte level within input data, as well as their capacity to recognize the substructure of tokens—an area where traditional tokenizer-based approaches often fall short. To assess these characteristics, we conducted supplementary experiments using benchmarks specifically designed to test both the robustness of models to noise in the input and their sensitivity to character-level information, across English and multiple languages, and incorporating elements such as digits and phonemes. The findings from these assessments are summarized in Table~\ref{tab:char_tasks}.

\paragraph{Noisy Data} To assess the resilience of different models, we introduce noised variants of the benchmark classification tasks outlined in Section~\ref{section:evals}, facilitating a direct comparison between tokenizer-based architectures and {\textsc BLT}. We implement five unique character-level perturbation techniques to modify the original text: (a) \textit{AntSpeak}: Alters the entire sequence into uppercase characters separated by spaces. (b) \textit{Drop}: Randomly deletes 10\% of characters from the text. (c) \textit{RandomCase}: Randomly assigns uppercase to half the characters and lowercase to the remaining half. (d) \textit{Repeat}: Randomly selects 20\% of the characters for repetition, up to four times. (e) \textit{UpperCase}: Converts all letters to uppercase. For evaluation, each noising method is applied separately to the prompt, the completion, or both, with average performance metrics reported. Table \ref{tab:char_tasks} presents the outcomes on noised HellaSwag~\citep{Eval_hellaswag}, demonstrating that {\textsc BLT} consistently exhibits greater robustness compared to tokenizer-based systems, achieving an average margin of 8 points over the counterpart trained on identical data, and even surpasses the Llama 3.1 model, which benefits from training on a substantially larger dataset.

\paragraph{Phonology - Grapheme-to-Phoneme (G2P)} We evaluate how effectively {\textsc BLT} translates sequences of graphemes (i.e., word-level characters) into their corresponding phonemic representations. The outcomes of the G2P experiment, conducted in a 5-shot configuration with Phonology Bench~\citep{suvarna2024phonologybench}, are displayed in Table \ref{tab:char_tasks}. Our findings indicate that {\textsc BLT} surpasses the performance of the Llama 3 1T tokenizer-based baseline in this setting.

\paragraph{CUTE}
To evaluate character-level comprehension, we test {\textsc BLT} on the CUTE benchmark~\citep{edman2024cute}, which consists of multiple tasks categorized into three main areas: compositional understanding, orthographic similarity recognition, and sequence manipulation capability. This benchmark represents a considerable challenge for most tokenizer-based systems, which often possess knowledge of how their tokens are spelled yet find it difficult to exploit this information for textual manipulation. As presented in Table~\ref{tab:char_tasks}, {\textsc BLT}-Entropy surpasses both BPE Llama 3 variants by over 25 points in this assessment. Notably, our approach excels in character manipulation tasks, achieving near-perfect accuracy (99.9\%) on both spelling-related evaluations. These significant gains, especially considering that {\textsc BLT} was trained on sixteen times less data than Llama 3.1, suggest that acquiring character-level knowledge is particularly challenging for BPE-based architectures. Figure \ref{fig:cute_examples} showcases selected cases where the Llama 3 tokenizer model falls short, yet {\textsc BLT} achieves successful outcomes. The only tasks where BPE models outperform are word deletion and insertion. While such word-level manipulations may not be inherently intuitive for a byte-level system, the performance gap is relatively minor, and constructing word operations from character-level capabilities may be more tractable than the reverse. We apply a consistent evaluation protocol across all experiments, utilizing the original Huggingface prompts. It is worth noting that BPE models could potentially see improved results with enhanced prompt engineering.

\paragraph{Low Resource Machine Translation} We assess the performance of {\textsc BLT} on both translation directions involving six major language families and a set of twenty-one low-resource languages, each representing different scripts, utilizing the FLORES-101 benchmark~\citep{goyal2022flores}. SentencePiece BLEU scores are presented in Table~\ref{tab:flores}. Our findings indicate that {\textsc BLT} achieves stronger results than a model leveraging the Llama 3 tokenizer, showing a 2-point overall improvement when translating into English and a 0.5-point gain when translating from English. For high-profile language pairs, {\textsc BLT} performs at least as well as, and occasionally better than, Llama 3. Notably, within many pairs from less-resourced language families, {\textsc BLT} demonstrates superior performance, highlighting the advantages of byte modeling for generalizing across the diverse and infrequent byte sequences characteristic of low-resource languages.

\subsection{Training {\textsc BLT} from Llama 3}
\label{sec:distillation}
We investigate an approach where {\textsc BLT} models can benefit from the capabilities of established pre-trained tokenizer-based models to enhance both training speed and performance. This is accomplished by initializing the global transformer parameters in {\textsc BLT} using those from a pre-trained Llama 3.1 model. Following this initialization, we proceed to fine-tune the global transformer weights with a learning rate set to one-tenth of that used for the local encoder and decoder components, applying this over the same number of training steps optimized for Llama 3. Comparative results, summarized in Table~\ref{tab:distill}, show that this {\textsc BLT} variant, initialized from Llama 3.1, achieves notable improvements over both the original Llama 3 and baseline {\textsc BLT}, under identical compute budgets. Notably, even when contrasted with the {\textsc BLT}-Entropy variant—trained with a much larger corpus of 1T tokens, as reported in Table~\ref{tab:evals}—the {\textsc BLT} initialized from Llama 3.1 demonstrates superior results on the MMLU benchmark. This indicates that the described strategy can substantially reduce the computational requirements needed for effective training.

This approach can be interpreted as converting models that rely on tokenizers into ones that do not, thereby transforming a pre-trained LLaMA 3.1 model into a {\textsc BLT} model. For a thorough evaluation, we report the results of the original LLaMA 3.1, which was trained on 15T tokens, in Table \ref{tab:distill}, and compare its performance to that of the {\textsc BLT} model generated from LLaMA 3. While our model shows only a modest decrease in accuracy on MMLU and HumanEval, there is a notably larger degradation across other benchmarks. These observations indicate that additional research is necessary to make better use of the pre-trained model's capabilities, with particular attention to refining data composition and optimizing other hyperparameters.

\section{Ablations and Discussion}
\label{section:ablations}
This section presents a series of ablation studies that rationalize the architectural design of {\textsc BLT} as well as the chosen patching strategy and associated hyper-parameters for the 8B parameter {\textsc BLT} model trained on the {\textsc BLT}-1T dataset.

\paragraph{Entropy Model Hyper-parameters} We systematically investigate how the size of the entropy model and the length of its context window influence scaling dynamics. To this end, we train byte-level entropy transformer models ranging from 1 million to 100 million parameters, each configured with context windows varying between 64 and 512 tokens. Scaling law relationships between bits-per-byte (bpb) and training FLOPs are visualized for our $400m$ and $1b$ {\textsc BLT} models, both trained on the Llama-2 dataset, in \autoref{fig:entropy_ablation}. Our experiments reveal that increasing either entropy model size or context window length leads to improved scaling, although the incremental benefits taper off beyond 50 million parameters.

\paragraph{Types of Patching}

We perform an ablation study on the four distinct patching strategies described in Section~\ref{section:patching}, namely: 1) Strided Patching employing strides of 4 and 6, 2) Patching at whitespace positions, 3) BPE Tokenizer patching utilizing the Llama 3 tokenizer, and 4) Entropy-driven patching leveraging a compact byte-level language model.

Although dynamic patching shortens the actual sequence lengths, we regulate sequence length across patching strategies to ensure that the context window remains consistent. Throughout training and inference, each model is exposed to an equivalent number of bytes per sequence on average, thereby eliminating confounding variables linked to modeling extended contexts. Figure~\ref{fig:scaling-compute-opt} presents the outcomes of these ablation studies. All alternative patching methods surpass static patching in performance, and space patching, in particular, shows results nearly matching those of dynamic entropy-based patching.

In \autoref{tab:evals_ablation}, we report benchmark results for {\textsc BLT} models trained on the Llama 2 dataset, evaluating the performance of tokenizer-based models alongside those utilizing space patching and entropy-based patching. Each model was trained for the optimal number of steps as established by~\citet{dubey2024llama}. While space patching represents a more straightforward method, avoiding the need for real-time entropy model computations during training, our analysis reveals that the improvements previously attributed to entropy-based patching in scaling experiments (see Section~\ref{section:scaling}) are also evident in downstream benchmark tasks.\footnote{The results for space patching are derived from earlier experiments conducted without cross-attention; however, comparable trends are also observed with cross-attention included.}

\paragraph{Cross-Attention}
\autoref{tab:crossattn_ablations_1b} presents our ablation study on the placement of cross-attention mechanisms within both the encoder and decoder components of {\textsc BLT}. Regarding the encoder, we investigate several query initialization strategies: (1) assigning an identical learned embedding to each global state, (2) applying a hash embedding to the bytes within each patch, and (3) deriving the query from a pooled representation of the encoder's hidden states corresponding to the patch bytes at the current encoder layer.

Our results indicate that incorporating cross-attention within the \textit{decoder} yields the highest effectiveness. When applied to the encoder, cross-attention confers a modest benefit, though this is apparent only when queries are initialized via pooling. Moreover, the use of cross-attention proves advantageous specifically for Common-Crawl data and shows even greater benefit when larger patch sizes are utilized.

\paragraph{n-gram Hash Embeddings} 

We evaluate the impact of varying the n-gram hash embedding vocabulary size across 0, 100K, 200K, and 400K settings, with the corresponding outcomes detailed in Table~\ref{tab:ngram_ablations_1b}. The results indicate that hash embeddings confer benefits across all tested domains, with especially pronounced improvements on Wikipedia and Github (0.04 bpb improvement compared to 0.01 bpb after 15k steps at the 8B scale). At the 8B parameter level, reducing the number of hashes from 500K to 300K leads to only a minor change in performance (~0.001 bpb at 15k steps). These findings suggest that hash embeddings play a crucial role in enabling {\textsc BLT} to achieve parity with tokenizer-based architectures; however, increases beyond 300K hashes yield progressively smaller gains. Furthermore, the enhancements from hash embeddings appear to work in conjunction with cross-attention, as improvements manifest across distinct datasets.

\paragraph{Local Model Hyperparameters} Table \ref{tab:local_layers} presents an ablation study on different choices for the layer counts in the local encoder and decoder. The results indicate that, particularly when utilizing hash n-gram embeddings, {\textsc BLT} performs optimally when the encoder is highly lightweight, consisting of only a single layer, while employing a more substantial decoder architecture.

\section{Related Work}
\label{section:related_work}

\textbf{Character-Level RNNs:} The task of Character Language Modeling has maintained significance in neural model research since its inception~\citep{sutskever2011generating,mikolov2012subword,graves2013generating}, largely due to its inherent ability to handle out-of-vocabulary tokens directly rather than relying on explicit back-off mechanisms. In their work, \cite{kim2016character} develop a model that exclusively ingests character-level information through convolutional and highway network layers, which subsequently provide input to LSTM-based RNNs. Their approach achieves parity with contemporary state-of-the-art RNN language models for English and demonstrates superior results for morphologically complex languages—highlighting an additional benefit of character-level LLMs. Similarly, \cite{kenter2018byte} address machine comprehension tasks utilizing byte-level LSTM architectures, surpassing word-level models in performance on languages with rich morphology such as Turkish and Russian. In another example, \cite{zhang2015character} employ convolutional networks at the character level for classification purposes, with results exceeding those of word-level approaches in specific tasks. Building upon these advancements, \citet{chung2019hierarchical} introduce hierarchical LSTM models equipped with boundary detectors at various tiers to uncover the latent hierarchical organization in text, achieving further enhancements for character-level language modeling. Additionally, ByteNet proposed by \cite{kalchbrenner2016neural} leverages convolutional neural networks applied at the character level instead of attention mechanisms for tasks such as machine translation.

\textbf{Character-Level Transformers:} The introduction of transformer architectures utilizing attention mechanisms~\citep{vaswani2017attention}, combined with advancements in subword tokenization~\citep{sennrich-etal-2016-neural}, has led to notable gains in neural language modeling and performance on standard benchmarks. Nevertheless, the use of word and sub-word representations introduces a specific inductive bias regarding the abstraction layer at which these models function. To bridge transformer approaches with early positive results from character-level language modeling, \citet{al2019character} employ very deep transformer architectures enhanced by auxiliary loss functions, allowing their transformer-based models to surpass earlier LSTM-driven character language models. Despite these advancements, they observed that character-level transformers still lagged noticeably behind models operating at the word level. Similarly, GPT-2~\citep{radfordgpt2} reported that, on expansive corpora such as the 1 Billion Word Benchmark, byte-level language models remained inferior to their word-level counterparts in terms of competitiveness.

\cite{Choe2019BridgingTG} showed that transformer-based byte-level LLMs can surpass subword-level LLMs with a similar parameter count; however, these models require significantly greater computational resources and substantially longer training times. Likewise, \cite{el2020characterbert} introduced CharFormer, a BERT variant that constructs word-level representations via convolutions over character embeddings. Their experiments in the medical domain produced improved outcomes, but the computational demands were substantially higher. The CANINE model developed by \cite{clark2022canine} consists of 150M parameters and directly processes character sequences. It leverages a deep transformer stack—analogous to the global model discussed here—and incorporates a local transformer in conjunction with strided convolutions for character downsampling. This architecture enabled CANINE to outperform its token-based encoder-only counterpart (mBERT) across multiple multilingual downstream tasks. ByT5~\citep{xue2022byt5} investigated byte-level encoder-decoder models that abstain from any patching mechanism. Their work showed improved robustness to input noise and competitive results compared to tokenizer-based models using four times less training data; yet, omitting patching led to highly expensive attention computations over all bytes, making these models very computationally intensive. Working at the byte level, rather than with subword units, further increases input sequence lengths, complicating efficient scaling of such models. More recently, \cite{wang2024mambabyte} leveraged the Mamba Architecture~\citep{gu2023mamba}, which is capable of maintaining a fixed-size memory state even with long context windows, to develop a byte-level Mamba model—again without patching. At a scale of 350M parameters, this model achieved lower bits-per-byte and outperformed comparable byte-level transformer models under a controlled FLOP budget across multiple datasets.

\textbf{Patching-based approaches:} Utilizing patching techniques has the potential to significantly reduce the computational cost (in terms of flops) associated with byte-level LLMs while maintaining comparable performance levels. Several studies have reported promising results on relatively modest model sizes and limited training data. For instance, \cite{nawrot-etal-2022-hierarchical} investigated static approaches to downsampling and upsampling through patching, introducing the hourglass transformer architecture, which demonstrated superior results over other byte-level methods at the 150M parameter level. Building on this, \cite{nawrot-etal-2023-efficient} introduced dynamic patching mechanisms, incorporating a boundary-predictor that is either trained end-to-end, guided through supervision with specific tokenizers, or constructed based on entropy—akin to {\textsc BLT}. Their findings indicated that these strategies can surpass the performance of contemporaneous vanilla transformers on language modeling benchmarks using a model with 40M parameters trained on approximately 400M tokens. In related work, \cite{lester2024training} explored the effect of training on sequences compressed via arithmetic coding, enabling compression rates that exceed those achievable by BPE. Through the implementation of the equal-info windows approach, they achieved improvements over byte-level baselines for language modeling, although their results still lagged behind those attained with subword baselines.

Our research is primarily motivated by MegaByte~\citep{yu2023megabyte}, a causal decoder-only LLM that processes data by applying a predetermined, fixed patching approach, concatenating representations to transform bytes into patches, and subsequently utilizing a local model on the decoder to reconstruct bytes from patches. The MegaByte architecture is shown to achieve performance comparable to tokenizer-based models at the 1B parameter scale on datasets totaling approximately 400B bytes. Throughout our experiments, we systematically compare with MegaByte and observe that the static patching strategy falls short of the most advanced, compute-efficient tokenizer-based models under controlled FLOP conditions. We then illustrate how {\textsc BLT} helps to close this performance gap. Similarly, \cite{slagle2024spacebyte} observe MegaByte's limitations and propose modifications such as applying patching based on whitespaces and similar characters, in addition to introducing a local encoder. Their results indicate benefits over tokenized transformer baselines in certain areas, including Github and arXiv, at the 1B parameter level under controlled compute. Our study extends this line of inquiry by presenting further experiments with this variant, concluding that additional architectural advances are essential for scaling byte-level models effectively to reach or surpass the performance of leading token-based architectures like Llama 3.

\section{Limitations and Future Work}

For the purposes of determining model architectures in this study, we follow the optimal training step counts identified for Llama~3~\citep{dubey2024llama}. It should be noted, however, that these scaling laws are derived based on BPE-level transformer models and may not yield ideal data-to-parameter ratios when applied to the {\textsc BLT} architecture. The derivation of scaling laws specifically tailored to {\textsc BLT} remains an avenue for future investigation and could reveal more advantageous scaling behavior for our model. Furthermore, most experiments here are limited to models up to 1B parameters, so it remains possible that the best architectural configurations may shift as the model size increases to 8B parameters or higher, potentially resulting in even stronger performance at larger scales.

Current transformer libraries and code repositories are optimized primarily for transformer models that utilize tokenizers. Although our work includes experiments that are theoretically matched in terms of floating point operations and incorporates some efficient implementations—like FlexAttention—to accommodate layers that diverge from the standard transformer setup, our implementations might still fall short of tokenizer-based models when it comes to actual wall-clock performance. Additional optimizations could further narrow this gap.

Whereas {\textsc BLT} employs an entropy model for patching that is trained independently, jointly optimizing the patching model within an end-to-end training pipeline could offer a promising avenue for subsequent research. In Section \ref{sec:distillation}, we detail preliminary experiments that provide early evidence supporting the effectiveness of converting large-scale tokenizer-based models—such as Llama 3, which has been trained on over 10T tokens—into byte-based models by reusing and fixing the pretrained global transformer weights. Continued investigation in this area could lead to approaches that not only preserve the advantages of bytefying, but also potentially surpass the performance of their tokenizer-based counterparts, all without the need to retrain the models from the beginning.

\section{Conclusion}
\label{section:conclusion}

In this work, we introduce the Byte Latent Transformer (\textbf{BLT}), an innovative architecture that moves beyond the traditional reliance on static tokenization schemes in large language models. {\textsc BLT} utilizes an adaptive and trainable approach to cluster bytes into variable-length patches, which allows the model to distribute computational effort according to the underlying complexity of the data. This mechanism results in notable gains in both computational efficiency and resilience. Through comprehensive scaling experiments, we find that {\textsc BLT} achieves parity with tokenization-centric models such as Llama 3 at model sizes up to 8B and datasets encompassing 4T bytes, while delivering as much as a 50\% decrease in inference FLOPs for only marginal compromises in accuracy. Additionally, {\textsc BLT} introduces an alternative scaling axis: it enables concurrent growth of both model and patch size while holding the inference computational budget constant—a feature particularly beneficial under typical deployment constraints. By processing raw bytes natively, {\textsc BLT} enhances the model’s capacity to manage rare or atypical data, improves robustness to input noise, and offers richer modeling of sub-word patterns. Collectively, these advantages highlight {\textsc BLT} as a compelling successor to conventional tokenization methodologies, advancing a more robust and resource-efficient foundation for adaptive language modeling.

\end{document}